
$\Gr K^{\bullet}$ (VI 5.5), on
peut supposer que $\underline{E}$ est un faisceau inversible. On a alors : $\underline{\check E}=\underline{E}^{-1}$ ; donc $\check x=x^{-1}$.

Les $c^i(\check x)$, $c^j(x)$ sont nuls pour $i,j\geq2$. Il suffit donc de montrer que $c^1(x^{-1})=-c^1(x)$. Or $c^1(x)$ est la classe dans $\Gr^kK^{\bullet}(X)$ de $x-1$, et $c^1(x^{-1})$ celle de $x^{-1}-1$. L'égalité :
\[
x^{-1}-1=1-x+(x-1)(1-x^{-1})
\]
donne :
\[
x^{-1}-1\equiv1-x\qquad\mathrm{mod.}\ \Fil^2K^{\bullet}(X);
\]
on en déduit le résultat.

\textbf{Théorème 3.6. (Riemann--Roch).} \uline{Soit} $f:X\to Y$ \uline{un morphisme projectif d'intersection complète} (1.1), \uline{et dont la dimension relative virtuelle} $d$ \uline{est constante} (1.9). \uline{On suppose que} $Y$ \uline{est quasi-compact}, \uline{et admet un faisceau ample}. \uline{Avec les notations de} 3.3, \uline{on a alors pour tout} $x\in K^{\bullet}(X)$ \uline{la relation}:
\[
\ch(f_*(x))=f_{\mathrm{gr}}\bigl(\ch(x)\cdot\Todd(\check T_f)\bigr),
\]
$\check T_f$ \uline{désignant l'élément dual de la classe du complexe cotangent relatif} $\underline{L}^{X/Y}_\bullet$ (2.5 \uline{et} 2.3).

En d'autres termes, le diagramme
\[
\begin{tikzcd}[column sep=huge,row sep=large]
K^{\bullet}(X) \arrow[r,"\ch\circ\widetilde c"] \arrow[d,"f_*"'] & \Gr K^{\bullet}(X)_{\mathbb Q} \arrow[d,"x\mapsto f_{\mathrm{gr}}(x\cdot\Todd(\check T_f))"]\\
K^{\bullet}(Y) \arrow[r,"\ch\circ\widetilde c"] & \Gr K^{\bullet}(Y)_{\mathbb Q}
\end{tikzcd}
\]
est commutatif.

